\section{Limitations, M2/M2.1 Split, and Conclusion}
\label{sec:conclusion}

\paragraph{Scope limitation.}
The positive flagship is a Lean Proof-layer theorem at the semantic axiom level.
It does not claim a result outside the finite \((\mathtt{Fin}\ n, \mathtt{Fin}\ m)\) family with \(2 \leq n\) and \(4 \leq m\); it does not claim a result for SWFs that fail \(\mathtt{UN}\) or \(\mathtt{MINLIB}\); and it does not claim a result that bypasses the polymorphic axiom definitions used at the base case.
The negative-side scope wall (Negative \#1) is a guarantee-boundary statement and an audit-policy statement; it is not (yet) a Lean theorem.

\paragraph{What \(\mathtt{\_hm}\) does not say.}
The hypothesis \(\mathtt{\_hm : 4 \leq m}\) is present in the theorem statement for symmetry with Sen's original formulation and to document the lift scope.
It is not consumed by the proof: the overlap case analysis avoids requiring a universal four-point completion in \(\mathtt{Fin\ m}\), since the two-overlap case yields an asymmetry contradiction, the one-overlap case yields a 3-cycle, and the disjoint case yields a 4-cycle, each assembled from alternatives already supplied by \(\mathtt{MINLIB}\).
We do not claim that \(\mathtt{UN \wedge MINLIB}\) implies \(\mathtt{4 \leq m}\).
The cardinality lemmas \(\mathtt{exists\_not\_mem\_of\_card\_lt}\) and \(\mathtt{four\_le\_card\_fin}\) remain in the lifting module as honest infrastructure documenting an avoided proof risk; they are not used by \(\mathtt{sen\_impossibility\_lifts}\).

\paragraph{M2/M2.1 split (forward pointer).}
M2 deliberately limits itself to the positive Proof-layer lift (Section~\ref{sec:positive}) and the Negative \#1 scope wall (Section~\ref{sec:negative}).
A separate Negative \#2 question---whether the \emph{representation} under which lever-level repairs are reported (bundled vs.\ split, clause-equivalent re-encodings) transfers canonically at family scale, i.e.\ carrying the M1.5 non-canonicity phenomenon to family scale---is real and methodologically required for the four-contract program, but it requires fixing the witness class and the generator interface before exploration scripts proliferate, and it does not block the present milestone.
It is therefore ratified as a separate milestone, \emph{M2.1}.
The decision is recorded in \code{docs/gates/m2/aim\_signoff.md}; the M2.1 reproduction record will live alongside that milestone, not in this paper.

\paragraph{What remains open.}
Two follow-ups are recorded explicitly.
First, promoting Proposition~\ref{prop:scope-wall} to a Lean theorem would require formalising the short-cycle encoding as a Lean-side relation and proving non-equivalence with full \(\mathtt{SocialAcyclic}\) at \(m = 5\); this is a non-blocking M2 follow-up.
Second, M2.1 will introduce a parametric witness-class interface and re-examine whether the M1.5 non-canonicity phenomenon recurs across family-scale clause-equivalent re-encodings; M2 makes no prediction here.

\paragraph{Conclusion.}
The contribution of M2 is a Transfer Contract for a finite social-choice impossibility witness, made along the layer-level boundary between the Lean Proof layer and the CNF Witness/Audit layer.
At the Proof layer, Sen's impossibility lifts unconditionally at the kernel-checked level to all finite cases with at least two voters and at least four alternatives, by direct polymorphic port of the base-case case-split; the hypothesis \(\mathtt{\_hm : 4 \leq m}\) is not consumed by the proof.
At the Witness/Audit layer, the lift does not automatically upgrade the sen24 CNF atlas into a family-level CNF certificate; the short-cycle encoding remains audited only at the base case, and a stronger family-level rationality encoding would have to be introduced and audited separately.
The methodological point is that ``the finite witness transferred'' is not a single statement: in an auditable-abstraction program, it is at least two statements, one per layer.
M2 makes the positive Proof-layer statement kernel-checked and the negative Witness/Audit-layer statement explicit, and defers the further representation-canonicity question to M2.1.
